\section{优先研究方向与实验设计}
\label{sec:future-directions}

以下方案由既有证据和模型候选共同提出，均为实验假设，而非已经验证的目标相配方。建议按“局部相图--位点占据/氧计量--晶体生长--类比取代”的顺序推进；前一阶段的相组成结果决定后一阶段是否值得开展。

\subsection{方向一：目标组成附近的局部相图}

第一优先级是在$\mathrm{Ba_5Y_{12}ZnSi_8O_{40}}$周围建立小步长组成网格。基础前驱体采用模型排名第一且化学上常规的BaCO$_3$、Y$_2$O$_3$、ZnO和SiO$_2$，目标化学计量对应$5{:}6{:}1{:}8$的摩尔比。本方向属于模型预测，待实验验证。建议先按无Zn固相研究的温区分段筛查：约1000\,$^\circ$C、12\,h进行初始反应，复磨后在1300\,$^\circ$C、24\,h退火，并依据PXRD结果决定是否重复；只有在仍存在明显未反应物且失重和坩埚相容性可控时，才增设1450--1600\,$^\circ$C的短时节点。温度和时间取自相近Ba--Y固相研究的筛查区间\cite{yamane2024microstructure,gulay2024navigation}，仅用于安排实验矩阵，不表示目标相必然在这些条件下形成。

组成轴至少应包含：(i)目标点附近Ba\slash Y比的小幅偏离；(ii) Zn含量的变化；(iii)与氧计量相关的无Zn参照相。每个样品均记录烧前和烧后质量、热历史和冷却方式，并以定量PXRD\slash Rietveld、SEM--EDS\slash EPMA以及必要的ICP分析给出相组成与实际元素比。若出现玻璃或多相边界，应在该区域加密取样，而不是仅保留最理想的衍射图谱。

\subsection{方向二：Zn--Y--氧计量耦合}

一个可检验的电荷平衡系列为
\[
\mathrm{Ba_5Y_{13-\mathit{x}}Zn_{\mathit{x}}[SiO_4]_8O_{8.5-\mathit{x}/2}},
\qquad 0\leq x\leq1.25,
\]
其中$x=0$对应已报道的无Zn四方参照相，$x=1$在元素计量上对应目标组成。建议以$x=0,0.25,0.50,0.75,1.00,1.25$为首轮离散点。本方向属于模型预测，待实验验证。前驱体仍采用BaCO$_3$、Y$_2$O$_3$和SiO$_2$，以ZnO调节Zn含量；该表达式仅为电荷平衡假设，不预设Zn的真实位点占据或氧空位分布。判断时应同时考察晶胞参数的连续性、相分数、实际Zn\slash Y比和位点占据精修；条件允许时，中子衍射或独立的氧含量与价态测量可进一步区分Zn取代、氧缺陷与两相混合。

该系列的科学价值高于单点复现：结构参数连续变化支持有限固溶；晶胞参数基本不变而两相比例变化则更接近两相区；仅在$x\approx1$出现新衍射峰则提示线化合物。三种结果均能约束目标相的形成机制。

\subsection{方向三：固相反应与高温溶液法晶体生长并行}

固相反应路线用于定位块体相区，高温溶液法路线用于获得可供SCXRD分析的晶体。本方向属于模型预测，待实验验证。高温溶液法路线可以Ba--Y--Si--O体系中已验证的K$_2$CO$_3$\slash MoO$_3$助熔体系为对照起点，前驱体采用方向一的BaCO$_3$、Y$_2$O$_3$、ZnO和SiO$_2$；在1100--1200\,$^\circ$C范围内设置最高温度，保温约3\,h，以1--2\,K\,h$^{-1}$缓冷至约900\,$^\circ$C，并系统改变营养物与助熔剂的比例以及Y$_2$O$_3$\slash SiO$_2$比\cite{kolitsch2009crystal,wierzbickawieczorek2017high}。这些数值来自无Zn同系列的晶体生长，必须与目标投料的固相反应路线并行比较，不能称为目标相配方。

每个晶体生长批次至少应报道液相组成、装填率、最高温度、保温时间、冷却速率、分离方法、晶体尺寸和产率；晶体以SCXRD和EDS验证，残余母液以PXRD和成分分析检查。若仅得到少量目标单晶而母液组成远离目标组成，应解释为液相中的选择性生长，而非块体相区。

\subsection{方向四：Mg\slash Co类比与模型前驱体}

本文调用本地两步无机逆合成模型，对三个元素式进行了Top-5组合预测。模型对Zn、Mg和Co目标的首选组合均保留BaCO$_3$、Y$_2$O$_3$和SiO$_2$，第四种原料分别为ZnO、MgO和Co$_3$O$_4$；三种首选组合在各自候选池中的概率分别为0.594、0.522和0.298。本方向属于模型预测，待实验验证。表~\ref{tab:mcp-experiments}\CJKecglue{}将模型输出转化为可检验的最小实验。模型仅依据元素计量提出前驱体组合，不用于判断硅酸根结构单元；正文始终保留目标文献所用的结构式。

\begin{table}[bp]
  \centering
  \caption{\textbf{两步无机逆合成模型给出的首选前驱体组合及其实验用途。}所有组合均为模型排序结果（未经化学路线验证）；工艺条件取自相近体系文献，需经相图实验复核。}
  \label{tab:mcp-experiments}
  \footnotesize
  \begin{tabular}{P{\dimexpr 0.1188\textwidth-2\tabcolsep\relax}P{\dimexpr 0.3100\textwidth-2\tabcolsep\relax}P{\dimexpr 0.1300\textwidth-2\tabcolsep\relax}P{\dimexpr 0.4412\textwidth-2\tabcolsep\relax}}
    \toprule
    \thead{假设目标} & \thead{模型首选前驱体} & \thead{候选池概率} & \thead{建议实验与判据}\\
    \midrule
    Zn目标 & BaCO$_3$ + Y$_2$O$_3$ + SiO$_2$ + ZnO & 0.594 & 作为局部相图的基础投料；先进行固相筛查，再将近单相组成转入助熔剂法晶体生长；以PXRD\slash Rietveld与SCXRD分别确定块体相区和结构归属\\
    \addlinespace[3pt]
    Mg类比 & BaCO$_3$ + Y$_2$O$_3$ + SiO$_2$ + MgO & 0.522 & 以Mg$^{2+}$作为同价尺寸对照；沿与Zn相同的组成/热历史矩阵筛查；仅当出现新相且实际组成连续变化时，才讨论同一位点取代\\
    \addlinespace[3pt]
    Co类比 & BaCO$_3$ + Y$_2$O$_3$ + SiO$_2$ + Co$_3$O$_4$ & 0.298 & 以空气气氛为首轮条件；测定实际组成和相分数；避免将第二相误判为结构取代\\
    \bottomrule
  \end{tabular}
\end{table}

模型给出的其余Zn路线包含额外的BaO、BaF$_2$、Ba(NO$_3$)$_2$或非标准的数据库标签，其概率均明显低于首选组合，不应并列推荐。额外Ba源会改变局部相图；BaF$_2$增加含氟挥发物和废物流的风险；非标准标签还可能反映训练数据的标准化问题。Mg路线中的MgCO$_3$可作为单独的前驱体形态对照，但不应与MgO同时加入；Co的多种氧化物前驱体不应视为三个等价方案。

\subsection{面向科学发现的实验组织}

每轮实验应以“最大化可区分结论”为标准，而不是仅以获得晶体为目标。信息量最大的组合为：(i)小步长组成网格配合定量相分析，建立局部相图；(ii)同一组成的固相反应与助熔剂法平行实验，区分平衡相区与液相中的选择性成核；(iii) Zn含量与氧计量联动系列，检验位点和缺陷假设。所有负结果同样予以保存，因为它们决定下一轮的组成选择，并约束“可合成”的边界。
